\section{Brain Connectivity Representation and Translation}

Learned representations of brain activity can be used not only for dimensionality reduction, but also to study how latent dimensions relate to functional brain organisation. 
Sanz Perl et al.\ used a variational autoencoder to obtain low-dimensional representations of parcellated fMRI activity~\cite{SanzPerl2025}. 
Individual latent dimensions were perturbed and decoded back into regional brain activity, allowing their spatial patterns to be related to large-scale functional networks. 
This showed that information represented within individual latent dimensions can be mapped back to interpretable patterns of functional brain organisation.

Latent representations can also provide a common space for translating between different representations of brain connectivity. 
Jamison et al.\ proposed Krakencoder, an encoder-decoder framework that maps different structural and functional connectomes into a shared latent representation~\cite{Jamison2025}. 
Separate encoders and decoders are used for different connectivity representations, allowing one representation to be encoded into the shared space and reconstructed as another. 
This demonstrates how a common latent representation can be combined with different decoders to reconstruct or translate between related representations of brain organisation.

Shared representations of structural and functional connectivity have also been investigated in relation to Alzheimer’s disease. 
Amodeo et al.\ proposed a graph variational autoencoder that combines structural connectivity with resting-state functional connectivity within a joint latent representation~\cite{Amodeo2022}. 
The learned representations were evaluated using data from the OASIS-3 dataset and contained information relevant to distinguishing AD-related populations. 
Together, these studies show that latent representations can preserve information about functional brain organisation, provide interpretable mappings back to brain networks, and support translation between different brain representations. 
These approaches provide a basis for investigating whether a shared representation can also be decoded into different disease-related functional brain states.
